\chapter{Requirements Analysis}
\label{chap:requirements}

\section{Introduction}

This chapter turns the context and objectives from Chapter~\ref{chap:framework} into a concrete requirements analysis. It starts with a short review of existing solutions, to position what CANvisionNative adds. It then lists the functional and non functional requirements, identifies the actors, and presents the use cases with a UML use case diagram and textual descriptions.

\section{Existing solutions analysis}

Tools an engineer might reach for today fall into three categories, none of which combine diagnostics, monitoring, replay, and explainability the way this project set out to.

\subsection{Category 1: CAN analyzers}

Commercial bench tools such as Vector CANoe and CANalyzer are the professional standard for observing \gls{can} traffic. They capture and display raw frames, decode them against a \gls{dbc} file, and let an engineer replay a recorded log with a rich, configurable interface. What they do not offer out of the box is runtime awareness: they are built around static monitoring and offline inspection, showing an engineer what is on the bus without judging whether a given frame is normal for the vehicle's current state. Turning them into something closer to that requires custom scripting on top of the base product.

\subsection{Category 2: IDS tools}

A separate line of published and research tooling targets the \gls{can} bus specifically for intrusion detection. These approaches generally fall into three families. \textbf{Signature-based detection} matches traffic against known attack patterns; it is precise for known attacks but cannot catch a new one. \textbf{Specification-based detection} encodes expected behaviour by hand (expected identifiers, frequency ranges, payload ranges); it does not need attack data, but it needs a manually written specification for every vehicle and does not generalize. \textbf{Anomaly-based detection} learns what normal traffic looks like from timing and/or payload statistics and flags deviations; it is the most broadly applicable of the three, and it is the detection family the anomaly-detection layer of CANvisionNative belongs to. In practice, published \gls{can} \gls{ids} tools tend to be rule-based or single-signal, are tuned toward catching attacks rather than explaining bus behaviour, and, being research-oriented, carry a high false-positive rate on real, unfiltered vehicle traffic outside their original benchmark.

\subsection{Category 3: Static viewers and diagnostic scanners}

A third category covers lighter-weight tools: open-source utilities such as candump and SavvyCAN, and \gls{obd} (SAE~J1979~\cite{sae_j1979}) or \gls{uds} (ISO~14229~\cite{iso14229}) diagnostic scanners. The open-source viewers display raw frames with no replay control, no anomaly awareness, and no unified dashboard. The diagnostic scanners take a different but equally narrow approach: they query specific \gls{ecu}s for known parameters (fault codes, sensor values) through a request/response protocol, so they only see what they explicitly ask for and have no notion of monitoring the bus continuously. Neither sub-type gives an engineer runtime intelligence about what the traffic as a whole is doing.

\subsection{Current limitations}

Across the three categories above, three limitations recur. First, most anomaly-based designs rely on a single detection signal (either timing or one machine-learning model), so an attack that does not disturb that one signal goes undetected. Second, a machine-learning model trained on one vehicle is rarely reused correctly across vehicles, because message layout, frequency, and normal-value ranges are vehicle-specific; in practice, systems that mention vehicle-specific models often treat correct model selection as an implementation detail rather than something to verify explicitly. Third, none of the tools reviewed above turn a raw anomaly score into a plain-language explanation an analyst can act on directly, and none combine diagnostics, monitoring, replay, and explanation in one workflow.

\subsection{Why CANvisionNative is different}

CANvisionNative is built as a diagnostic and validation platform first: its purpose is to give a BMS engineer one place to import, replay, monitor, and understand \gls{can} traffic, rather than to compete as a specialized detector against the IDS tools reviewed above. Around that diagnostic core, it combines several capabilities that, individually, exist in the categories above but never together in one tool: continuous monitoring instead of only offline inspection, log replay across multiple formats instead of a single proprietary one, machine-learning-based anomaly detection instead of a bare display of raw frames, and a natural-language explanation layer instead of a numeric score with no context.

On the detection side specifically, it fuses four independent signals (timing, payload, machine learning, temporal persistence) into one score, so an attack or fault missed by one layer can still be caught by another, and it selects a trained model per vehicle profile at run time rather than relying on a single generic model; this exact selection step was one of the six defects verified during validation (Chapter~\ref{chap:validation}). Security is a real strength of this design, since the same anomaly awareness that flags a sensor fault will also flag a malicious frame, but it is a strength that follows from the diagnostic core, not the reason the platform exists.

\section{Functional requirements}

\begin{itemize}
    \item \textbf{FR1 --- Import a recorded log.} The system shall let the user import a \gls{can} log file in \gls{csv}, \gls{mf4}, Vector \gls{can} (.asc), candump (.log), or PCAN text (.txt) format, and convert it to a canonical frame stream (timestamp, identifier, payload).
    \item \textbf{FR2 --- Select a vehicle profile.} The system shall let the user select the vehicle that produced the log (or the closest matching profile), so that the correct \gls{dbc} file and the correct trained model are used.
    \item \textbf{FR3 --- Replay the log.} The system shall replay the imported frames through the detection pipeline, in order, at a user-selectable speed (0.5$\times$ up to a maximum speed), and report progress.
    \item \textbf{FR4 --- Detect anomalies.} For every frame, the system shall compute a fusion anomaly score from four detection layers (timing, payload, machine learning, temporal persistence) and raise an alert when the score exceeds the configured threshold.
    \item \textbf{FR5 --- Classify severity.} Each alert shall be labelled with a severity level (LOW, MEDIUM, HIGH, CRITICAL) derived from the fusion score.
    \item \textbf{FR6 --- Explain an alert.} The system shall let the user request a natural-language explanation of any alert, including the dominant detection layer, a plain-language description of the likely cause, and a recommended action.
    \item \textbf{FR7 --- Export results.} The system shall let the user export the alert list as a \gls{csv} file and as a formatted PDF report.
    \item \textbf{FR8 --- Live simulation.} The system shall be able to generate a simulated live \gls{can} stream, including a simulated attack burst, for testing the detection pipeline without a physical vehicle.
    \item \textbf{FR9 --- Live hardware capture (software interface).} The system's software shall expose a live hardware capture interface, connecting to a physical \gls{can} adapter (ELM327 or SocketCAN) and feeding received frames into the same detection pipeline used for replay. This requirement covers the software interface only: it was implemented and is code-complete, but it was not experimentally validated against a physical vehicle within this project (Chapter~\ref{chap:validation}, Limitations).
    \item \textbf{FR10 --- Record a session.} The system shall be able to record a live or simulated session to a log file, optionally injecting a labelled attack, for later replay or retraining.
    \item \textbf{FR11 --- Configure the system.} The system shall let the user configure the alert threshold, the \gls{llm} provider key, and the hardware interface settings, and persist this configuration between sessions.
\end{itemize}

\section{Non-functional requirements}

\begin{itemize}
    \item \textbf{NFR1 --- Performance.} The replay engine shall process a one-hour, 97{,}741-frame recording in a few minutes, not hours, on a standard desktop machine. Chapter~\ref{chap:validation} reports a measured throughput of about 286 frames per second after the corrections described there.
    \item \textbf{NFR2 --- Reliability.} A detection defect must be provable and fixable with concrete evidence, not by blindly tuning thresholds. This principle governed the whole validation phase in Chapter~\ref{chap:validation}.
    \item \textbf{NFR3 --- Usability.} An analyst without machine-learning expertise shall be able to import a log, read the alert list, and understand the explanation of an alert without consulting external documentation.
    \item \textbf{NFR4 --- Portability of input formats.} Every supported log format shall enter the detection pipeline through the same canonical frame representation, so that detection behaviour does not depend on which format was used to record the data.
    \item \textbf{NFR5 --- Extensibility.} Adding a new vehicle profile shall only require adding a \gls{dbc} file and a trained model directory, without changing the detection engine code.
    \item \textbf{NFR6 --- Graceful degradation.} If the \gls{llm} provider is unavailable, the system shall still produce a rule-based explanation instead of failing the request.
    \item \textbf{NFR7 --- Platform.} The desktop client shall run on Windows 10/11 using .NET Framework 4.8 and \gls{wpf}; the backend shall run as a local Python process exposing a \gls{rest} \gls{api} on localhost. The full software and hardware requirements are listed in Table~\ref{tab:app-requirements} and Table~\ref{tab:app-hardware} (Appendix~\ref{app:config}).
\end{itemize}

\section{Actors}

\begin{itemize}
    \item \textbf{Analyst (primary user).} Imports logs, starts replays, reviews alerts, requests explanations, exports reports, configures the system.
    \item \textbf{Vehicle / CAN bus (external system).} Source of live frames when the hardware capture mode is used.
    \item \textbf{LLM provider (external system).} Supplies natural language text for the explanation layer (Gemini, or a local Ollama model as fallback).
    \item \textbf{CANvisionNative backend (system under design).} Performs decoding, feature extraction, scoring, and explanation generation, and exposes them through the \gls{rest} \gls{api}.
\end{itemize}



\section{Use cases}

Figure~\ref{fig:usecase} shows the use case diagram for the Analyst actor, the two external systems, and the main use cases implemented by the system.

\begin{figure}[H]
\centering
\includegraphics[width=\textwidth]{usecase.png}
\caption{Use case diagram for CANvisionNative.}
\label{fig:usecase}
\end{figure}

Import log, Select vehicle profile, Review alerts, Export report, Configure settings, Start live simulation, and Connect hardware are all initiated directly by the Analyst. Detect anomaly is included by Replay log, Start live simulation, and Connect hardware, since all three input modes feed the same detection pipeline. Request explanation extends Review alerts, since it is an optional action available on any reviewed alert, and it involves the external LLM provider. Table~\ref{tab:usecase-table} gives a short textual description of the eight use cases most representative of the system's behaviour.

\begin{table}[H]
\centering
\caption{Textual description of the main use cases.}
\label{tab:usecase-table}
\begin{tabularx}{\textwidth}{|p{3.2cm}|X|}
\hline
\textbf{Use case} & \textbf{Description} \\
\hline
Import log & The analyst selects a log file; the system detects its format and converts it to the canonical frame stream (FR1). \\
\hline
Select vehicle profile & The analyst chooses the vehicle that produced the log, so the correct \gls{dbc} file and trained model are used (FR2). \\
\hline
Replay log & The system replays the canonical frame stream through the detection pipeline at a user-selectable speed, including the Detect anomaly use case (FR3, FR4). \\
\hline
Detect anomaly & For every frame, the system computes the four-layer fusion score and raises an alert above the configured threshold (FR4, FR5). Shared by replay, live simulation, and hardware capture. \\
\hline
Review alerts & The analyst browses the alert list, filters by severity, and inspects the layer-attribution breakdown of a selected alert. \\
\hline
Request explanation & The analyst asks for a natural-language explanation of a selected alert; the system builds context and queries the LLM provider chain, falling back to a rule-based template if needed (FR6). \\
\hline
Export report & The analyst exports the current alert list as \gls{csv} or as a formatted PDF report (FR7). \\
\hline
Configure settings & The analyst sets the alert threshold, the LLM provider key, and the hardware interface parameters, persisted between sessions (FR11). \\
\hline
\end{tabularx}
\end{table}

\subsection{Activity diagram --- import and replay}

Figure~\ref{fig:activity-replay} shows the activity flow from importing a log to publishing its alerts to the client, corresponding to the Import log, Select vehicle profile, Replay log, and Detect anomaly use cases together.

\begin{figure}[H]
\centering
\includegraphics[width=0.5\textwidth]{activity_replay.png}
\caption{Activity diagram: import and replay a log.}
\label{fig:activity-replay}
\end{figure}

\subsection{Sequence diagram --- request an AI explanation}

Figure~\ref{fig:seq-explain} shows the sequence of calls triggered by the Request explanation use case, including the fallback behaviour when the primary LLM provider is unavailable.

\begin{figure}[H]
\centering
\includegraphics[width=0.85\textwidth]{sequence_explain.png}
\caption{Sequence diagram: requesting an AI explanation for an alert.}
\label{fig:seq-explain}
\end{figure}

\section{Risk analysis}

Table~\ref{tab:risks} lists the main risks identified during the project and how each was mitigated.

\begin{table}[H]
\centering
\caption{Risk analysis.}
\label{tab:risks}
\begin{tabular}{|p{3.3cm}|p{2.7cm}|p{1.7cm}|p{6.5cm}|}
\hline
\textbf{Risk} & \textbf{Likelihood} & \textbf{Impact} & \textbf{Mitigation} \\
\hline
No compatible electric vehicle available for hardware validation & High (materialized) & Medium & Design and implement the hardware layer against the same interface as replay and simulation, so it can be validated later without changing the detection pipeline (Chapter~\ref{chap:design}). \\
\hline
LLM provider unavailable or unauthenticated & Medium & Low & Provider chain with automatic fallback: Gemini $\rightarrow$ local Ollama model $\rightarrow$ rule-based template. \\
\hline
Vehicle-specific model not selected correctly & Medium (materialized) & High & Root-cause investigation with direct evidence rather than threshold tuning; full fix and validation documented in Chapter~\ref{chap:validation}. \\
\hline
False positives eroding analyst trust & Medium & High & Multi-layer fusion with a temporal-persistence term instead of a single-frame score; before/after alert-rate measurement in Chapter~\ref{chap:validation}. \\
\hline
Different log formats behaving inconsistently & Medium (materialized) & Medium & Single canonical parsing entry point shared by every input path, verified in Chapter~\ref{chap:validation}. \\
\hline
Explanation shown for the wrong alert (stale cache) & Low (materialized) & High & Root-caused to two cache-lifecycle gaps and fixed at both the completion and the start of every replay, verified with a seeded cross-dataset test (Chapter~\ref{chap:validation}). \\
\hline
\end{tabular}
\end{table}

\section{Planning}

The project was carried out over the full six-month internship, through a sequence of milestones tracked internally as M1 to M8, followed by a dedicated final validation phase, as introduced in the development methodology of Chapter~\ref{chap:framework}. Figure~\ref{fig:gantt} presents this planning as a Gantt chart. \todofill{Confirm exact calendar start/end dates and per-milestone week ranges; the chart below preserves the original relative ordering and overlaps, stretched to the real six-month (26-week) duration.}

\begin{figure}[H]
\centering
\resizebox{\textwidth}{!}{%
\begin{ganttchart}[
    hgrid, vgrid,
    bar/.append style={fill=cyan!60},
    bar height=0.6,
    bar label font=\small,
    y unit chart=0.7cm,
    x unit=0.55cm
]{1}{26}
\gantttitle{Project weeks (6 months)}{26} \\
\gantttitlelist{1,...,26}{1} \\
\ganttbar{M1 --- Live foundation}{1}{3} \\
\ganttbar{M3 --- Multi-state ML training}{3}{8} \\
\ganttbar{M4 --- AI explanation engine}{6}{11} \\
\ganttbar{M5 --- Real-time graphs}{8}{12} \\
\ganttbar{M6 --- Session recording}{11}{14} \\
\ganttbar{M8 --- Offline / live decode pipeline}{10}{16} \\
\ganttbar{M7 --- Integration}{14}{17} \\
\ganttbar{M2 --- Hardware layer (code, EV validation pending)}{14}{17} \\
\ganttbar{Root-cause validation (real data)}{17}{21} \\
\ganttbar{Engineering correction \& re-validation}{20}{24} \\
\ganttbar{Report writing}{23}{26}
\end{ganttchart}%
}
\caption{Project planning (Gantt chart, in weeks, over the six-month internship).}
\label{fig:gantt}
\end{figure}

Milestone M2 (hardware / OBD2 validation) is shown as code-complete but is explicitly marked as blocked, not on the hardware adapter itself, but on the absence of a compatible electric vehicle to validate against during the project window. It is not part of the validated scope in Chapter~\ref{chap:validation}, which is why it is not part of the risk-mitigated critical path.

\section{Requirements traceability}

Table~\ref{tab:traceability} traces each functional requirement to the design and implementation section that realizes it and, where applicable, the validation evidence that confirms it works as intended.

\begin{table}[H]
\centering
\caption{Requirements traceability: from requirement to implementation to validation.}
\label{tab:traceability}
\begin{tabularx}{\textwidth}{|p{1.3cm}|X|p{3.4cm}|}
\hline
\textbf{Req.} & \textbf{Implemented in} & \textbf{Validated by} \\
\hline
FR1 & Canonical parser, Section~\ref{sec:impl-parser} & Fix 4, Chapter~\ref{chap:validation} \\
\hline
FR2 & Vehicle profile system, Section~\ref{sec:impl-vehicle} & Fix 1, Chapter~\ref{chap:validation} \\
\hline
FR3 & Replay engine, Section~\ref{sec:impl-replay} & Before/after throughput, Table~\ref{tab:before-after} \\
\hline
FR4 & Detection pipeline, Section~\ref{sec:design-pipeline} & Root-cause analysis, Table~\ref{tab:rootcause} \\
\hline
FR5 & Fusion scoring, Section~\ref{sec:design-pipeline} & Severity distribution, Figure~\ref{fig:chart-severity} \\
\hline
FR6 & AI explanation module, Section~\ref{sec:design-ai} & Fix 6, Section~\ref{sec:validation-ai} \\
\hline
FR7 & Export, Section on Export, Chapter~\ref{chap:impl} & Export validation, Chapter~\ref{chap:validation} \\
\hline
FR8 & Live simulation (hardware layer design) & Interactive UI audit, Chapter~\ref{chap:validation} \\
\hline
FR9 & Hardware layer software interface, Section on Hardware, Chapter~\ref{chap:impl} & Code-complete; physical vehicle test pending, no compatible EV available (Limitations) \\
\hline
FR10 & Session recording (backend/hardware/) & Interactive UI audit, Chapter~\ref{chap:validation} \\
\hline
FR11 & Settings, Section on Settings, Chapter~\ref{chap:impl} & Interactive UI audit, Chapter~\ref{chap:validation} \\
\hline
\end{tabularx}
\end{table}

Every functional requirement is either backed by direct, reproducible validation evidence or explicitly marked as pending a physical test, consistent with the honest treatment of the hardware capture limitation throughout this report.

\section{Conclusion}

This chapter reviewed existing \gls{can} tools and intrusion detection approaches, positioned CANvisionNative against their limitations, and defined the functional and non functional requirements, actors, and use cases that follow from this analysis, together with the project's risk analysis and planning. The next chapter presents the system design built to satisfy these requirements.
